\begin{table}[t!]
\caption{Results of agreements between two human annotation results and between human and model evaluation results. }
\vspace{-0.075in}
\label{tab:agreement}
\small
\centering
\resizebox{0.475\textwidth}{!}{
\renewcommand{\arraystretch}{1.0}
\renewcommand{\tabcolsep}{1.0mm}
\begin{tabular}{llccc}
\toprule

\textbf{Categories} & \textbf{Metrics} & \textbf{Problem} & \textbf{Method} & \textbf{Experiment} \\

\midrule
\midrule

\multirowcell{2}[-0.0ex][l]{\textbf{Human and Human}} 

& Scoring & 0.83 & 0.76 & 0.67 \\

& Pairwise & 0.62 & 0.62 & 0.41 \\

\midrule

\multirowcell{2}[-0.0ex][l]{\textbf{Human and Model}} 

& Scoring & 0.64 & 0.58 & 0.49 \\

& Pairwise & 0.71 & 0.62 & 0.52 \\

\bottomrule

\end{tabular}
}
\vspace{-0.025in}
\end{table}


\section{Experimental Results and Analyses}
\label{sec:results}

\paragraph{Main Results}
Our main results on scoring with human and model-based evaluations are provided in Figure~\ref{fig:main}. These demonstrate that our full ResearchAgent outperforms all baselines by large margins on every metric across problems, methods, and experiment designs (constituting the complete research ideas). Particularly, the full ResearchAgent augmented with relevant entities exhibits strong gains on metrics related to creativity (such as Originality for problems and Innovativeness for methods) since entities may offer novel concepts and views that may not be observable in the group of citation-based papers alone. In addition, the results of pairwise comparisons between models with both human and model-based evaluations -- shown in Figure~\ref{fig:pairwise} -- demonstrate that the full ResearchAgent shows the highest win ratio over its baselines.  

\paragraph{Analysis on Inter-Annotator Agreements}
To validate the quality and reliability of human annotations, we measure the inter-annotator agreements, where 20\% of the generated ideas are evaluated by two human judges, and report results in Table~\ref{tab:agreement}. Specifically, for the scoring, we first rank scores from each annotator and measure Spearman's correlation coefficient~\cite{spearman} between the ranked scores of two annotators. For the pairwise comparison between two judges, we measure Cohen’s kappa coefficient~\cite{cohen}. Table~\ref{tab:agreement} shows that the inter-annotator agreement is high, confirming the reliability of our assessments about the quality of generated research ideas. Also, while agreement scores for experimental designs are slightly lower than other aspects, this does not necessarily indicate a shortcoming in the quality of experimental designs produced by ResearchAgent, as demonstrated in Figures~\ref{fig:main} and~\ref{fig:pairwise}. Instead, we view this as the inherent subjectivity and variability in how such designs are perceived and evaluated by different annotators (i.e., the nature of the variability itself makes achieving high agreement challenging).

\begin{figure}[t!]
    \centering
    \includegraphics[width=0.985\columnwidth]{Figures/iteration_fig.pdf}
    \vspace{-0.08in}
    \caption{Results with varying the number of refinement steps.}
    \label{fig:iteration}
    \vspace{-0.07in}
\end{figure}
\begin{table}[t!]
\caption{Results of ablation study on references and entities.}
\vspace{-0.1in}
\label{tab:ablation}
\small
\centering
\resizebox{0.475\textwidth}{!}{
\renewcommand{\arraystretch}{0.8}
\begin{tabular}{lccc}
\toprule

\textbf{Methods} & \textbf{Problem} & \textbf{Method} & \textbf{Experiment} \\

\midrule
\midrule

ResearchAgent & \textbf{4.52} & \textbf{4.28} & \textbf{4.18} \\

\noalign{\vskip 0.25ex}\cdashline{1-4}\noalign{\vskip 0.75ex}

- w/o Entities & 4.35 & 4.13 & 4.02 \\

- w/ Random Entities & 4.41 & 4.19 & 4.13 \\

\noalign{\vskip 0.25ex}\cdashline{1-4}\noalign{\vskip 0.75ex}

- w/o References & 4.26 & 4.08 & 3.97 \\

- w/ Random References & 4.35 & 4.16 & 4.02 \\

\noalign{\vskip 0.25ex}\cdashline{1-4}\noalign{\vskip 0.75ex}

- w/o Entities \& References & 4.20 & 4.03 & 3.92 \\

\bottomrule

\end{tabular}
}
\vspace{-0.1in}
\end{table}


\paragraph{Analysis on Human-Model Agreements} Similar to what we did for the aforementioned inter-annotator agreements, we measure agreements between human-based and model-based evaluations, to ensure the reliability of model-based evaluations. As shown in Table~\ref{tab:agreement}, we further confirm that agreements between humans and models are high, indicating that model-based evaluations are a reasonable proxy to judge research idea generation. 

\paragraph{Analysis of Refinement Steps}
To see the effectiveness of iterative refinements of research ideas with ReviewingAgents, in Figure~\ref{fig:iteration}, we report the averaged scores on the generated ideas as a function of refinement steps. We first observe initial improvements in the quality of generated ideas with increased refinement steps. Yet, the performance becomes saturated after three iterations, which may indicate diminishing returns for subsequent iteration steps, which aligns with the pattern observed in agent-based refinement work~\cite{multiagent}.


\paragraph{Ablation on Knowledge Sources}
Recall that the full ResearchAgent is augmented with two different knowledge sources, namely relevant references and entities. To see their individual contribution, we perform an ablation study by either excluding one of the knowledge sources or replacing it with random elements. As shown in Table~\ref{tab:ablation}, each knowledge source contributes to performance improvement, and the relevant references are especially helpful. We also note that providing random elements is more helpful than providing no elements at all; we hypothesize that this may be due to the LLM's capability to filter out noise while still gaining incidental value from random inputs.

\begin{figure}[t!]
    \centering
    \includegraphics[width=0.975\columnwidth]{Figures/distribution_fig.pdf}
    \vspace{-0.1in}
    \caption{Distributions of model-based evaluation results with and without the human-induced score criteria alignment (middle and right), as well as human evaluation results (left).}
    \label{fig:distribution}
    \vspace{-0.05in}
\end{figure}
\begin{figure}[t!]
    \centering
    \includegraphics[width=0.975\columnwidth]{Figures/citation_fig.pdf}
    \vspace{-0.1in}
    \caption{Results with bucketing papers based on citations.}
    \label{fig:citation}
    \vspace{-0.075in}
\end{figure}

\paragraph{Analysis on Human Alignment for Evaluation}
Recall that to align judgments from model-based evaluations with actual human preferences, we generated the evaluation criteria based on human evaluation results and used them as the criteria for model-based evaluations. Figure~\ref{fig:distribution} demonstrates the efficacy of this strategy, presenting the score distribution of human evaluation compared with the distributions of model-based evaluations with and without human alignment. We find that the score distribution of model-based evaluations without alignment is skewed and different from the score distribution of human judgments. Meanwhile, after aligning the model-based evaluations with human-induced score criteria, the calibrated distribution more closely resembles the distribution of humans.

\paragraph{Correlation on Citation Counts}
We further investigate whether a high-impact paper (when used as a core paper) leads to high-quality research ideas. To measure this, we categorize papers by their citation count (as a proxy for impact), and visualize the average score of each bucket (with model-based evaluations) in Figure~\ref{fig:citation}. We find that ideas from high-impact papers tend to be of higher quality, likely due to their ability to identify research gaps, propose feasible methods, and connect with other works. Additionally, based on the paper distribution (See Figure~\ref{fig:discipline}) and for the ease of manual quality check, evaluation criteria for model-based evaluations are induced mainly with computer science papers. To see whether those criteria are applicable to diverse fields, we also compare a correlation between scores of computer science papers and all papers in Figure~\ref{fig:citation}. From this, we observe that the scores increase when the citation increases for both domains, which may support the generalizability of human-preference-induced evaluation criteria.

\begin{table}[t!]
\caption{Comparisons of ResearchAgent with hypothesis generation methods~\cite{CLBD, tomato}.}
\vspace{-0.1in}
\label{tab:hypothesis}
\small
\centering
\resizebox{0.475\textwidth}{!}{
\renewcommand{\arraystretch}{0.8}
\renewcommand{\tabcolsep}{0.75mm}
\begin{tabular}{lccccc}
\toprule

\textbf{Methods} & \textbf{Clarity} & \textbf{Relevance} & \textbf{Originality} & \textbf{Feasibility} & \textbf{Significance} \\

\midrule
\midrule

SciMON & 4.04 & 4.37 & 4.56 & 3.98 & 4.15 \\

Hypothesis Proposer & 3.97 & 4.14 & 4.07 & 4.01 & 4.11 \\

\noalign{\vskip 0.25ex}\cdashline{1-6}\noalign{\vskip 0.75ex}

ResearchAgent & \textbf{4.11} & \textbf{4.88} & \textbf{4.77} & \textbf{4.05} & \textbf{4.81} \\

\bottomrule

\end{tabular}
}
\vspace{-0.1in}
\end{table}


\paragraph{Comparisons to Hypothesis Generation}
Recall that existing methods for hypothesis generation focus on predicting links between variables or generating hypotheses based on these links, which differs from our experimental setup of generating open-ended research ideas (problems, methods, and experiments). Nevertheless, to understand how the quality of the generated research ideas from prior works~\cite{CLBD, tomato} differs from our ResearchAgent, we perform comparisons. As shown in Table~\ref{tab:hypothesis}, we observe that ResearchAgent is capable of generating superior research hypotheses, due to the utilization of broad and deep knowledge across domains as well as the iterative review and refinement procedures.

\paragraph{Analysis using Different LLMs}
To assess how ResearchAgent's performance changes with different LLMs, we conduct an auxiliary analysis with Llama3, Mixtral, Qwen1.5, and GPT-3.5~\cite{Qwen, Mixtral}, as shown in Table~\ref{tab:gpt35}. These results show a significant performance drop with less capable models. Moreover, the performance differences between the Naive ResearchAgent without knowledge augmentation and the full ResearchAgent become marginal (for Mixtral and GPT-3.5), which indicates that they might struggle with capturing complex concepts between scientific papers. This can likely be attributed to the emergent abilities of LLMs for complex reasoning (but not in smaller LMs)~\cite{Wei2022EmergentAO}, although other subtle issues may also be contributing factors.

\paragraph{Qualitative Analysis}
We provide qualitative results on generated research ideas in Table~\ref{tab:examples}. One representative example in the last row highlights the advantage of entity-centric knowledge augmentation, where two entities (such as Drosophila Genetic Reference Panel and CRISPR) retrieved from the entity-centric knowledge store enable the generation of a novel research idea: bridging genetic variability and CRISPR applications. This exemplifies how external entity-based knowledge uncovers non-trivial relations between scientific concepts.


\begin{table}[t!]
\caption{Results with different, open and proprietary LLMs.}
\vspace{-0.1in}
\label{tab:gpt35}
\small
\centering
\resizebox{0.475\textwidth}{!}{
\renewcommand{\arraystretch}{1.25}
% \renewcommand{\tabcolsep}{2.0mm}
\begin{tabular}{llccc}
\toprule

\textbf{LLMs} & \textbf{Models} & \textbf{Problem} & \textbf{Method} & \textbf{Experiment} \\

\midrule
\midrule

\multirowcell{2}[-0.0ex][l]{\textbf{GPT-4.0}} 

& Naive ResearchAgent & 4.20 & 4.03 & 3.92 \\

& ResearchAgent (Ours) & 4.52 & 4.28 & 4.18 \\

\midrule

\multirowcell{2}[-0.0ex][l]{\textbf{GPT-3.5}} 

& Naive ResearchAgent & 3.56 & 3.56 & 3.63 \\

& ResearchAgent (Ours) & 3.58 & 3.58 & 3.60 \\

\midrule

\multirowcell{2}[-0.0ex][l]{\textbf{Llama3 (8B)}} 

& Naive ResearchAgent & 3.76 & 3.69 & 3.54 \\

& ResearchAgent (Ours) & 4.18 & 4.03 & 3.95 \\

\midrule

\multirowcell{2}[-0.0ex][l]{\textbf{Mixtral (8x7B)}} 

& Naive ResearchAgent & 3.31 & 3.27 & 3.20 \\

& ResearchAgent (Ours) & 3.28 & 3.35 & 3.31 \\

\midrule

\multirowcell{2}[-0.0ex][l]{\textbf{Qwen1.5 (32B)}} 

& Naive ResearchAgent & 3.64 & 3.74 & 3.66 \\

& ResearchAgent (Ours) & 4.02 & 3.97 & 3.94 \\

\bottomrule

\end{tabular}
}
\vspace{-0.1in}
\end{table}
